% ============================================================================
% Fichier cible : ARIEL_LOGMINER_MEMOIRE_FINAL/chapters/limites_consolidees.tex
%   (§5.8 « Limites du travail », paragraphe « Grandeurs non mesurées. »,
%   PDF p. 113)
% Opération : AJOUT de deux phrases à la fin du paragraphe, après « [...] n'ont
%   fait l'objet d'aucune étude auprès d'analystes. »
% Auteur : Agent 2 (statut des revendications), chantier 2026-11-05.
% Motifs : absence de scénario hétérogène conjoint (L3, point de
%   non-recevabilité) ; priorisation, traitements périodiques et micro-lots
%   annoncés mais jamais mesurés (L2 §2.1 et faiblesse 7). La règle du corpus
%   (HANDOFF §1, règle 3) place les limites ici, une seule fois.
% Sources : §4.12 (corpus synthétique ; incidents enchaînant plusieurs
%   catégories d'événements) ; aucune métrique de priorisation ni campagne
%   périodique ou par micro-lots dans chapters/*.tex (micro-lots décrits au
%   §2.23.2 seulement).
% Arbitrage : ce fichier touche un périmètre partagé (limites) ; le
%   coordinateur décide de son intégration.
% ============================================================================

Aucune expérience n'analyse conjointement un même incident à partir de journaux système et de flux réseau réels : les sources sont évaluées séparément, et le seul regroupement d'événements de catégories différentes porte sur le corpus synthétique de corrélation. La priorisation des anomalies, les traitements périodiques et l'exécution par micro-lots, prévus par la conception, ne sont pas évalués.
